\begin{filecontents*}{refs.bib}
@inproceedings{wang2021tent,
  title={Tent: Fully Test-Time Adaptation by Entropy Minimization},
  author={Wang, Dequan and Shelhamer, Evan and Liu, Shaoteng and Olshausen, Bruno and Darrell, Trevor},
  booktitle={International Conference on Learning Representations},
  year={2021},
  url={https://arxiv.org/abs/2006.10726}
}
@inproceedings{wang2022cotta,
  title={Continual Test-Time Domain Adaptation},
  author={Wang, Qin and Fink, Olga and Van Gool, Luc and Dai, Dengxin},
  booktitle={IEEE/CVF Conference on Computer Vision and Pattern Recognition},
  year={2022},
  url={https://arxiv.org/abs/2203.13591}
}
@inproceedings{niu2022eata,
  title={Efficient Test-Time Model Adaptation without Forgetting},
  author={Niu, Shuaicheng and Wang, Yuqi and Wu, Jiaxiang and Yu, Lequan and Xu, Min},
  booktitle={International Conference on Machine Learning},
  year={2022},
  url={https://proceedings.mlr.press/v162/niu22a.html}
}
@inproceedings{iwasawa2021t3a,
  title={Test-Time Classifier Adjustment Module for Model-Agnostic Domain Generalization},
  author={Iwasawa, Yusuke and Matsuo, Yutaka},
  booktitle={Advances in Neural Information Processing Systems},
  year={2021},
  url={https://papers.nips.cc/paper/2021/hash/1415fe9fea0fa1e45dddcff5682239a0-Abstract.html}
}
@inproceedings{jang2023tast,
  title={Test-time Adaptation via Self-Training with Nearest Neighbor Information},
  author={Jang, Minguk and Chung, Sae-Young and Chung, Hye Won},
  booktitle={International Conference on Learning Representations},
  year={2023},
  url={https://arxiv.org/abs/2207.10792}
}
@inproceedings{niu2024foa,
  title={Test-Time Model Adaptation with Only Forward Passes},
  author={Niu, Shuaicheng and Miao, Chunyan and Chen, Guohao and Wu, Pengcheng and Zhao, Peilin},
  booktitle={International Conference on Machine Learning},
  year={2024},
  url={https://dblp.org/rec/conf/icml/NiuMCWZ24}
}
@article{xiao2024beyond,
  title={Beyond Model Adaptation at Test Time: A Survey},
  author={Xiao, Zehao and Snoek, Cees G. M.},
  journal={arXiv preprint arXiv:2411.03687},
  year={2024},
  url={https://arxiv.org/abs/2411.03687}
}
@inproceedings{carlini2022lira,
  title={Membership Inference Attacks from First Principles},
  author={Carlini, Nicholas and Chien, Steve and Nasr, Milad and Song, Shuang and Terzis, Andreas and Tramer, Florian},
  booktitle={IEEE Symposium on Security and Privacy},
  year={2022},
  url={https://arxiv.org/abs/2112.03570}
}
@inproceedings{wang2025reliability,
  title={Membership Inference Attacks as Privacy Tools: Reliability, Disparity and Ensemble},
  author={Wang, Zhiqi and Zhang, Chengyu and Chen, Yuetian and Baracaldo, Nathalie and Kadhe, Swanand and Yu, Lei},
  booktitle={Proceedings of the 2025 ACM SIGSAC Conference on Computer and Communications Security},
  year={2025},
  url={https://arxiv.org/abs/2506.13972}
}
@inproceedings{zheng2024ldpicl,
  title={Locally Differentially Private In-Context Learning},
  author={Zheng, Chunyan and Sun, Keke and Zhao, Wenhao and Zhou, Haibo and Jiang, Lixin and Song, Shaoyang and Zhou, Chunlai},
  booktitle={Proceedings of the Joint International Conference on Computational Linguistics, Language Resources and Evaluation},
  year={2024},
  url={https://arxiv.org/abs/2405.04032}
}
@inproceedings{wen2024iclmia,
  title={Membership Inference Attacks Against In-Context Learning},
  author={Wen, Rui and Li, Zheng and Backes, Michael and Zhang, Yang},
  booktitle={Proceedings of the ACM SIGSAC Conference on Computer and Communications Security},
  year={2024},
  url={https://arxiv.org/abs/2409.01380}
}
@article{duan2024privacyicl,
  title={On the Privacy Risk of In-context Learning},
  author={Duan, Haonan and Dziedzic, Adam and Yaghini, Mohammad and Papernot, Nicolas and Boenisch, Franziska},
  journal={arXiv preprint arXiv:2411.10512},
  year={2024},
  url={https://arxiv.org/abs/2411.10512}
}
@article{levin2025promptmia,
  title={Has My System Prompt Been Used? Large Language Model Prompt Membership Inference},
  author={Levin, Roman and Cherepanova, Valeriia and Hans, Abhimanyu and Schwarzschild, Avi and Goldstein, Tom},
  journal={arXiv preprint arXiv:2502.09974},
  year={2025},
  url={https://arxiv.org/abs/2502.09974}
}
@article{duan2023flocks,
  title={Flocks of Stochastic Parrots: Differentially Private Prompt Learning for Large Language Models},
  author={Duan, Haonan and Dziedzic, Adam and Papernot, Nicolas and Boenisch, Franziska},
  journal={arXiv preprint arXiv:2305.15594},
  year={2023},
  url={https://arxiv.org/abs/2305.15594}
}
\end{filecontents*}

\documentclass{article}

\IfFileExists{neurips_2024.sty}{
  \usepackage[final]{neurips_2024}
}{
  \usepackage[margin=1in]{geometry}
  \usepackage{microtype}
  \usepackage[numbers,sort&compress]{natbib}
}

\usepackage[T1]{fontenc}
\usepackage[utf8]{inputenc}
\usepackage{times}
\usepackage{hyperref}
\usepackage{url}
\usepackage{booktabs}
\usepackage{amsmath}
\usepackage{amssymb}
\usepackage{graphicx}
\usepackage{xcolor}
\usepackage{multirow}

\title{Who Was in the Recent Window?\\A Rigorous Audit of Online Test-Time Adaptation Privacy}

\author{Anonymous Authors}

\newcommand{\recentlabel}{m_t^H}
\newcommand{\ft}{f_t}
\newcommand{\fsrc}{f_{\mathrm{src}}}
\newcommand{\auc}{\mathrm{ROC\mbox{-}AUC}}

\begin{document}

\maketitle

\begin{abstract}
Online test-time adaptation (TTA) updates a deployed model on unlabeled user queries, creating a deployment-time privacy surface absent from frozen inference. We study one exact question: given black-box access to the deployed state $\ft$ at audit time $t$, can an attacker infer whether a candidate sample appeared in the recent adaptation window $[t-H+1,\ldots,t]$? We instantiate this audit on CIFAR-10-C with two corruptions, three seeds, 5{,}000-example streams, and three methods: source-only inference, Tent, and T3A. The benchmark combines output-only logistic auditors with unmatched, matched, and nearest-neighbor-matched control pools that remove shortcuts based on class, confidence, stream position, and local feature similarity. Tent delivers large utility gains, improving mean stream accuracy from $0.2794$ to $0.7292$ on gaussian noise and from $0.6571$ to $0.8484$ on motion blur. Privacy leakage after strict controls is much weaker: the strongest pooled NN-matched $\auc$ is $0.5589$ for T3A on motion blur at $H{=}8$, while Tent stays near chance on gaussian noise and only reaches $0.5441$ on motion blur at $H{=}32$. Low-FPR reporting is mostly underpowered under the final matched-pool sizes; at $H{=}32$, the largest pooled TPR@1\% FPR is $0.0188$. Overall, the benchmark is reproducible and useful, but our confirmatory evidence for strong black-box recent-window leakage is conditional rather than conclusive.
\end{abstract}

\section{Introduction}

Online test-time adaptation (TTA) updates deployed model state on unlabeled target inputs to maintain accuracy under distribution shift \citep{wang2021tent,wang2022cotta,niu2022eata}. That makes deployment-time privacy different from standard training-time membership inference. If the deployed state changes in response to recent queries, then those recent queries may become inferable from the adapted model, even when the original training checkpoint remains fixed.

This paper studies one narrow question:
\begin{quote}
Given black-box access to the deployed state $\ft$ at audit time $t$, can an attacker infer whether a candidate sample $x^\star$ appeared in the recent adaptation window $[t-H+1,\ldots,t]$?
\end{quote}
The formulation is deliberately specific. The protected event is indexed by time, so the audited model is also time-indexed: we audit $\ft$, not a final post-stream model. We do not propose a new TTA algorithm. Instead, we provide a benchmark and protocol for privacy auditing in online TTA, using matched-control design choices motivated by modern low-FPR membership inference practice \citep{carlini2022lira,wang2025reliability}.

The initial hypothesis behind the project was that at least one standard online TTA method would leak recent-window membership in a way that survives strict matching. The resulting evidence is weaker than that hypothesis. Tent produces large utility gains but little convincing black-box leakage after strict controls. T3A exhibits modest above-chance $\auc$ values, but the signal is inconsistent across horizons and is accompanied by elevated source-only and far-past controls. The paper therefore becomes a rigorous conditional result: under this benchmark, strong black-box recent-window leakage is not yet cleanly established.

Our contributions are:
\begin{itemize}
  \item We formalize recent-window membership for online TTA as a time-indexed audit problem on the deployed state $\ft$ and instantiate a reproducible black-box benchmark with matched and nearest-neighbor-matched controls.
  \item We run a complete CIFAR-10-C study over $18$ full streams (three methods, two corruptions, three seeds), with $153$ audit times per 5{,}000-example stream and horizons $H \in \{1,8,32\}$.
  \item We show a sharp utility--privacy mismatch: Tent improves mean stream accuracy by $0.4498$ on gaussian noise and $0.1913$ on motion blur relative to source-only inference, yet strict-control privacy leakage remains weak, with the strongest pooled NN-matched $\auc$ only $0.5589$ and the largest pooled TPR@1\% FPR only $0.0188$.
  \item We identify a practical evaluation bottleneck: balanced matched pools remove obvious shortcuts, but they leave most low-FPR estimates underpowered at $H{=}1$ and $H{=}8$, so matched and NN-matched $\auc$ become the primary confirmatory endpoints in this benchmark.
\end{itemize}

Section~\ref{sec:related} reviews prior work, Section~\ref{sec:method} defines the threat model and audit protocol, Section~\ref{sec:experiments} presents the experimental results, Section~\ref{sec:discussion} discusses limitations, and Section~\ref{sec:conclusion} concludes.

\section{Related Work}
\label{sec:related}

\paragraph{Test-time adaptation.}
Tent introduced entropy-minimization-based fully test-time adaptation via online updates to batch-normalization affine parameters \citep{wang2021tent}. Subsequent work broadened the design space with continual adaptation \citep{wang2022cotta}, selective and forgetting-aware updates \citep{niu2022eata}, classifier-state adaptation \citep{iwasawa2021t3a}, nearest-neighbor self-training \citep{jang2023tast}, and forward-only adaptation \citep{niu2024foa}. Recent surveys emphasize that test-time adaptation is now a deployment paradigm rather than one algorithm family \citep{xiao2024beyond}. These papers focus on utility under shift, not on whether recent deployment queries are inferable from the current deployed state.

\paragraph{Membership inference and privacy auditing.}
Low-FPR evaluation is central in modern privacy auditing. LiRA showed that privacy claims can change substantially when evaluated at strict false-positive rates \citep{carlini2022lira}, and subsequent work argued that privacy metrics should be analyzed for reliability and disparity rather than treated as single-run point estimates \citep{wang2025reliability}. Our work inherits that perspective, but the protected set is not the training set. It is a transient recent deployment window generated by online adaptation.

\paragraph{Adjacent deployment-time privacy.}
Recent prompt and in-context learning work shows that inference-time context can itself leak private information. This includes locally private in-context learning \citep{zheng2024ldpicl}, membership inference against in-context examples \citep{wen2024iclmia}, broader privacy analyses of in-context learning \citep{duan2024privacyicl}, prompt-membership inference for system prompts \citep{levin2025promptmia}, and differentially private prompt learning \citep{duan2023flocks}. These papers are the closest adjacent literature to our setting because they study transient deployment-time state. However, they focus on explicit prompt or context membership. Our question instead concerns implicit state updates induced by online TTA.

Unlike prior work on TTA, our benchmark audits recent deployment queries. Unlike prior work on in-context privacy, our benchmark does not rely on explicit prompt membership. It targets whether online model adaptation leaves a detectable trace of recent users in the deployed state itself.

\section{Method}
\label{sec:method}

\subsection{Threat Model}

Let $\fsrc$ be a source model and let an online TTA method process an unlabeled target stream $x_1,x_2,\ldots$, producing deployed states $\ft$ for times $t=1,2,\ldots$. For a horizon $H$, define recent-window membership
\begin{equation}
\recentlabel(x^\star) =
\begin{cases}
1 & \text{if } x^\star \in \{x_{t-H+1},\ldots,x_t\},\\
0 & \text{otherwise.}
\end{cases}
\end{equation}
The audit problem is: from black-box outputs of $\ft(x^\star)$, infer $\recentlabel(x^\star)$.

The primary attacker is black-box. It can query $\ft$ on a candidate image and observe class probabilities. It does not know the true label, internal weights, gradients, optimizer state, or adaptation statistics.

\subsection{Attack Features and Auditor}

For each query, we extract five output-space features from the adapted probabilities:
\begin{itemize}
  \item maximum predicted probability,
  \item entropy,
  \item top-1 minus top-2 margin,
  \item Gini concentration,
  \item $\ell_2$ norm of log-probabilities.
\end{itemize}
We train a class-balanced logistic-regression auditor on these five features. Audit times are split contiguously: the earliest $60\%$ are used for training, the next $20\%$ for validation and threshold selection, and the final $20\%$ for held-out testing. This avoids leakage across adjacent deployed states.

\subsection{Matched-Control Benchmark Construction}

The benchmark is created offline but mirrors the online threat model. At each audit time $t$, positives are the current recent-window members. For every positive anchor, we construct several negative pools:
\begin{itemize}
  \item \textbf{Unmatched}: random non-members, capped at $32$ negatives per positive.
  \item \textbf{Matched}: exact matching on corruption, severity, true class, source predicted class, source-confidence bin, and stream-time bucket, capped at $5$ negatives per positive.
  \item \textbf{NN-matched}: the matched pool filtered further by cosine nearest neighbors in frozen source-feature space, keeping the top $k=5$ neighbors.
  \item \textbf{Far-past}: matched negatives that appeared earlier in the stream but at least $4H$ positions before time $t$.
\end{itemize}

True labels are used only for offline benchmark curation, never by the black-box attacker. The confidence bins are three equal-frequency bins computed from source-model confidence, and stream position is coarsened into eight time buckets across the 5{,}000-example stream.

T3A is selective, so for T3A we additionally match on whether a candidate crosses the adaptation-acceptance threshold in the queried deployed state. This prevents ``was admitted for adaptation'' from becoming a shortcut.

\subsection{TTA Methods Under Audit}

\paragraph{Source-only.}
The source-only baseline performs frozen inference with the same backbone and preprocessing as the adaptive methods. It is the primary null model: in principle, recent-window membership should be difficult to infer because the deployed state does not change.

\paragraph{Tent.}
Following \citet{wang2021tent}, Tent updates only batch-normalization affine parameters by minimizing prediction entropy on each incoming batch. In our implementation, non-BN parameters are frozen, batch-normalization running statistics are disabled, and the optimizer is Adam with learning rate $10^{-3}$.

\paragraph{T3A.}
Following \citet{iwasawa2021t3a}, T3A keeps the backbone frozen and adapts classifier supports online. We initialize supports from the source classifier weights, accept new supports only when the maximum class probability is at least $0.15$, and keep the $32$ lowest-entropy supports per class. This produces a lighter-weight mechanism than gradient-based Tent, but it still changes the deployed decision state over time.

\subsection{Evaluation Metrics}

We report three classes of metrics.

\paragraph{Utility.}
For each run we measure mean stream top-1 accuracy over all $5{,}000$ stream elements, tail accuracy over the final $500$ elements, and gain over source-only inference.

\paragraph{Privacy.}
The main privacy metric is $\auc$. We also report TPR@1\% FPR and TPR@0.1\% FPR when the validation and test pools each have enough negatives. The benchmark enforces minimum sizes of $10{,}000$ negatives for 1\% FPR and $50{,}000$ negatives for 0.1\% FPR.

\paragraph{Uncertainty.}
We use stratified bootstrap confidence intervals over audit times with $1{,}000$ resamples. Bootstrap groups are defined at the audit-time level rather than the candidate-row level to preserve within-state dependence.

\section{Experiments}
\label{sec:experiments}

\subsection{Setup}

The mandatory study uses CIFAR-10-C with two corruptions, \texttt{gaussian\_noise} and \texttt{motion\_blur}, both at severity $5$. We use a single RobustBench CIFAR-10 corruption model (\texttt{Standard}) for all methods and all seeds. Each run replays one stream of length $5{,}000$ sampled without replacement from the $10{,}000$ corrupted test examples. The batch size is $64$, the adaptation chunk size is $32$, and audit times are every $32$ positions from $t=128$ to $t=4992$, giving $153$ audit times per full stream. We evaluate horizons $H \in \{1,8,32\}$ and seeds $\{0,1,2\}$.

The environment smoke test reports Python $3.12.7$, PyTorch $2.10.0{+}\mathrm{cu128}$, one NVIDIA RTX A6000 GPU, and four CPU cores. Mean runtime per 5{,}000-example stream is approximately $88.0$--$89.2$ seconds for source-only, $357.1$--$369.6$ seconds for Tent, and $340.5$--$352.9$ seconds for T3A.

\begin{figure}[t]
\centering
\includegraphics[width=0.95\linewidth]{figures/threat_model_schematic.png}
\caption{Threat model overview. A fixed source model $\fsrc$ is deployed with online TTA on an unlabeled target stream, producing intermediate deployed states $\ft$. At audit time $t$, the attacker queries $\ft$ on a candidate sample and tries to infer whether that sample was present in the recent adaptation window $[t-H+1,\ldots,t]$.}
\label{fig:threat}
\end{figure}

\subsection{Main Utility Results}

Table~\ref{tab:utility} shows that the utility behavior is highly method dependent. Tent substantially improves accuracy on both corruptions, while T3A is essentially neutral on motion blur and harmful on gaussian noise. This matters for privacy interpretation: any privacy cost should be judged against the utility gain actually obtained by the adaptation method.

\begin{table}[t]
\caption{Utility on CIFAR-10-C severity-5 streams. Results are mean $\pm$ sample standard deviation over three seeds. Gain is relative to the source-only baseline for the same corruption. Higher is better for all columns except gain can be negative.}
\label{tab:utility}
\centering
\resizebox{\linewidth}{!}{
\begin{tabular}{llccc}
\toprule
Corruption & Method & Mean Stream Acc. $\uparrow$ & Tail Acc. $\uparrow$ & Gain vs. Source \\
\midrule
\multirow{3}{*}{Gaussian noise}
& Source-only & 0.2794 $\pm$ 0.0026 & 0.2713 $\pm$ 0.0090 & 0.0000 $\pm$ 0.0000 \\
& Tent & \textbf{0.7292 $\pm$ 0.0080} & \textbf{0.7420 $\pm$ 0.0020} & \textbf{0.4498 $\pm$ 0.0057} \\
& T3A & 0.2517 $\pm$ 0.0273 & 0.2727 $\pm$ 0.0378 & -0.0277 $\pm$ 0.0262 \\
\midrule
\multirow{3}{*}{Motion blur}
& Source-only & 0.6571 $\pm$ 0.0081 & 0.6627 $\pm$ 0.0190 & 0.0000 $\pm$ 0.0000 \\
& Tent & \textbf{0.8484 $\pm$ 0.0083} & \textbf{0.8467 $\pm$ 0.0231} & \textbf{0.1913 $\pm$ 0.0038} \\
& T3A & 0.6587 $\pm$ 0.0095 & 0.6780 $\pm$ 0.0092 & 0.0015 $\pm$ 0.0156 \\
\bottomrule
\end{tabular}
}
\end{table}

\subsection{Main Privacy Results}

Table~\ref{tab:privacy} reports pooled matched and NN-matched $\auc$ values. Three observations stand out.

First, strict-control privacy leakage is much weaker than unmatched privacy would suggest. Tent is nearly at chance on gaussian noise for $H=8$ and $H=32$ and only modestly above chance on motion blur at $H=32$.

Second, T3A has the strongest above-chance results, with NN-matched pooled $\auc$ values of $0.5526$ on gaussian noise at $H=8$, $0.5569$ on gaussian noise at $H=32$, $0.5589$ on motion blur at $H=8$, and $0.5513$ on motion blur at $H=32$.

Third, the recency story is not monotone. Our original hypothesis expected the strongest signal at small horizons, but the most stable values often appear at $H=8$ or $H=32$ rather than $H=1$. That weakens the case that the auditor is isolating strictly recent exposure.

\begin{table}[t]
\caption{Pooled privacy results on CIFAR-10-C. Entries are matched and NN-matched pooled ROC-AUC on the held-out test blocks. Best NN-matched value within each corruption is shown in bold. Higher is worse for privacy.}
\label{tab:privacy}
\centering
\resizebox{\linewidth}{!}{
\begin{tabular}{llcccccc}
\toprule
Corruption & Method & M $H{=}1$ & NN $H{=}1$ & M $H{=}8$ & NN $H{=}8$ & M $H{=}32$ & NN $H{=}32$ \\
\midrule
\multirow{3}{*}{Gaussian noise}
& Source-only & 0.5400 & 0.5367 & 0.4995 & 0.5146 & 0.5166 & 0.5309 \\
& Tent & \textbf{0.5875} & \textbf{0.5782} & 0.5097 & 0.5073 & 0.5061 & 0.5054 \\
& T3A & 0.5424 & 0.5563 & \textbf{0.5335} & \textbf{0.5526} & \textbf{0.5393} & \textbf{0.5569} \\
\midrule
\multirow{3}{*}{Motion blur}
& Source-only & 0.4479 & 0.4641 & 0.5268 & 0.5431 & 0.5303 & 0.5442 \\
& Tent & \textbf{0.5445} & \textbf{0.5428} & 0.5197 & 0.5190 & 0.5280 & 0.5441 \\
& T3A & 0.5137 & 0.5253 & \textbf{0.5556} & \textbf{0.5589} & \textbf{0.5402} & \textbf{0.5513} \\
\bottomrule
\end{tabular}
}
\end{table}

The confidence intervals tell the same story. For example, T3A on motion blur at $H{=}8$ reaches pooled NN-matched $\auc = 0.5589$ with 95\% CI $[0.5391, 0.5780]$, whereas Tent on gaussian noise at $H{=}8$ is $0.5073$ with 95\% CI $[0.4735, 0.5420]$. However, source-only baselines are not perfectly flat either: source-only motion blur reaches NN-matched $\auc = 0.5442$ at $H{=}32$. This residual source-only signal is a warning that benchmark construction, not only adaptation, can still leave detectable structure.

\begin{figure}[t]
\centering
\includegraphics[width=0.95\linewidth]{figures/control_strength_h8.png}
\caption{Control-strength comparison at horizon $H{=}8$. Each bar is pooled ROC-AUC for unmatched (U), matched (M), and nearest-neighbor-matched (NN) evaluation. T3A preserves modest signal under stronger controls; Tent is close to chance once matching is applied.}
\label{fig:control}
\end{figure}

\subsection{Low-FPR Results}

Low-FPR metrics are mostly unavailable under the final matched-pool sizes. The benchmark requires at least $10{,}000$ negatives in both validation and test pools for TPR@1\% FPR, which fails for all $H{=}1$ and $H{=}8$ matched and NN-matched conditions. Only $H{=}32$ has enough negatives for pooled TPR@1\% FPR, and no condition reaches the $50{,}000$-negative threshold for TPR@0.1\% FPR.

Table~\ref{tab:lowfpr} reports the resulting $H{=}32$ values. Even there, the effect sizes are small. The largest pooled NN-matched TPR@1\% FPR is $0.0188$ for T3A on gaussian noise, and source-only baselines already reach $0.0151$ and $0.0168$. We therefore do not interpret the low-FPR results as strong evidence of privacy leakage.

\begin{table}[t]
\caption{Pooled TPR@1\% FPR at horizon $H{=}32$. Entries are reported only when both validation and test pools contain at least $10{,}000$ negatives. Higher is worse for privacy.}
\label{tab:lowfpr}
\centering
\begin{tabular}{llcc}
\toprule
Corruption & Method & Matched & NN-matched \\
\midrule
\multirow{3}{*}{Gaussian noise}
& Source-only & 0.0108 & 0.0151 \\
& Tent & 0.0097 & 0.0114 \\
& T3A & 0.0141 & \textbf{0.0188} \\
\midrule
\multirow{3}{*}{Motion blur}
& Source-only & 0.0134 & 0.0168 \\
& Tent & 0.0134 & \textbf{0.0181} \\
& T3A & --- & 0.0128 \\
\bottomrule
\end{tabular}
\end{table}

\subsection{Ablations and Controls}

Figure~\ref{fig:nulls} shows the null controls at $H{=}8$. For Tent, shuffled-membership and far-past controls remain near the matched recent-window values, which supports the interpretation that there is little clean privacy signal after strict control. For T3A, matched recent-window performance is above shuffled nulls, but far-past controls remain competitive, especially on motion blur. That again makes the evidence conditional rather than decisive.

Table~\ref{tab:ablation} summarizes the ablations aggregated over both corruptions and all three seeds. Two trends are clear. First, the multi-feature logistic auditor consistently outperforms the best scalar score, often by more than $0.05$--$0.10$ $\auc$ points. Second, relaxing individual controls changes the measured privacy only modestly. For T3A at $H{=}8$, the full NN-matched average $\auc$ is $0.5557$, while removing source-prediction matching reduces it to $0.5292$ and removing time-bucket matching reduces it to $0.5245$. For Tent, all variants remain near chance.

\begin{table}[t]
\caption{Ablation summary aggregated over both corruptions and all three seeds. ``Full'' is the black-box NN-matched ROC-AUC. The remaining columns remove one control from benchmark construction or replace the logistic auditor with the best scalar score chosen on seed 0. Higher values indicate worse privacy.}
\label{tab:ablation}
\centering
\resizebox{\linewidth}{!}{
\begin{tabular}{llccccccc}
\toprule
Method & $H$ & Full & $-$TrueCls & $-$Time & $-$SrcPred & $-$Conf & $-$NN & Best Scalar \\
\midrule
Tent & 8 & 0.5132 & 0.5322 & 0.5152 & 0.5121 & 0.5167 & 0.5190 & 0.4776 \\
Tent & 32 & 0.5247 & 0.5129 & 0.5217 & 0.5262 & 0.5325 & 0.5200 & 0.4695 \\
T3A & 8 & \textbf{0.5557} & 0.5355 & 0.5245 & 0.5292 & 0.5469 & 0.5452 & 0.4927 \\
T3A & 32 & \textbf{0.5541} & 0.5370 & 0.5256 & 0.5328 & 0.5496 & 0.5411 & 0.4888 \\
\bottomrule
\end{tabular}
}
\end{table}

\begin{figure}[t]
\centering
\includegraphics[width=0.95\linewidth]{figures/recency_decay.png}
\caption{NN-matched ROC-AUC versus horizon. The expected monotone recency decay is not observed cleanly. Tent is near chance on gaussian noise across all horizons, whereas T3A stays modestly above chance across $H \in \{8,32\}$ on both corruptions.}
\label{fig:recency}
\end{figure}

\begin{figure}[t]
\centering
\includegraphics[width=0.95\linewidth]{figures/null_controls_h8.png}
\caption{Null controls at $H{=}8$. For each method, we compare matched recent-window auditing against a far-past control and a shuffled-membership null. The gap between recent and null conditions is small for Tent and only moderate for T3A, especially once far-past controls are included.}
\label{fig:nulls}
\end{figure}

\section{Discussion and Limitations}
\label{sec:discussion}

The strongest empirical finding in this study is not a privacy break. It is the gap between utility and privacy evidence. Tent produces clear utility benefits, but once we enforce matched and NN-matched controls, its privacy signal is usually near chance. T3A shows more consistent above-chance $\auc$, but the effect is modest and does not clearly track recency. That matters because the motivating claim was specifically about \emph{recent-window} membership.

Several limitations follow directly from the current workspace.

\paragraph{Low-FPR underpowering.}
The final balanced matched pools are too small for most low-FPR estimates. This is not a reporting choice; it is enforced by the minimum-negative thresholds built into the analysis pipeline. As a result, matched and NN-matched $\auc$ become the primary confirmatory endpoints, and the strongest privacy claims that can be made are limited.

\paragraph{Residual baseline structure.}
Source-only and far-past controls are not perfectly flat. For example, source-only motion blur reaches NN-matched $\auc = 0.5442$ at $H{=}32$. This suggests that stream position, image hardness, or source-feature locality are not fully neutralized by the current benchmark design.

\paragraph{Scope.}
The study is intentionally narrow: CIFAR-10-C only, two corruptions, one severity, one backbone, and three methods. We do not claim transfer to other datasets, larger models, or modalities. We also do not evaluate stronger gray-box or white-box attackers.

\paragraph{Benchmark rather than proof of non-leakage.}
Near-chance black-box results do not prove that online TTA is private. They only show that, under this attack interface and benchmark construction, strong recent-window leakage was not cleanly demonstrated.

The broader lesson is methodological. Deployment-time privacy audits need the same discipline as modern training-time membership inference: exact threat models, strict matching, null controls, and enough negatives for low-FPR evaluation. Without those ingredients, it is too easy to confuse incidental structure with privacy leakage.

\section{Conclusion}
\label{sec:conclusion}

We presented a benchmark for auditing whether a candidate sample appeared in the recent adaptation window of an online TTA model state $\ft$. The benchmark fixes the threat model, enforces matched and nearest-neighbor-matched controls, and evaluates black-box auditors on CIFAR-10-C streams.

The resulting evidence is mixed and mostly conservative. Tent improves mean stream accuracy from $0.2794$ to $0.7292$ on gaussian noise and from $0.6571$ to $0.8484$ on motion blur, but its strict-control privacy signal is weak. T3A reaches the highest pooled NN-matched $\auc$ values, up to $0.5589$, yet low-FPR results remain small and source-only controls are not flat. We therefore view the current paper as a rigorous conditional audit result rather than a strong positive privacy attack.

Future work should scale the benchmark to larger pools, stronger attackers, and more realistic deployment settings so that low-FPR evaluation becomes routine rather than optional. That is the main path toward deciding whether recent deployment queries in online adaptation are measurably exposed in practice.

\bibliographystyle{plainnat}
\bibliography{refs}

\end{document}
